\begin{theorem}[最大子式综合征的 Hankel-正规形与唯一递推]\label{thm:xi-hankel-maximal-minor-syndrome-normal-form-uniqueness}
\leanverified{paper\_xi\_hankel\_maximal\_minor\_syndrome\_normal\_form\_uniqueness}
设 \(k\) 为域，\((a_n)_{n\ge 0}\subset k\) 的 Hankel 秩为 \(d<\infty\)，并假设 Hankel 主块
\[
H_0:=\bigl(a_{i+j}\bigr)_{0\le i,j\le d-1}\in M_d(k)
\]
可逆。
令扩展 Hankel 块
\[
\widehat H:=\bigl(a_{i+j}\bigr)_{\substack{0\le i\le d-1\\ 0\le j\le d}}\in M_{d,d+1}(k).
\]
对 \(0\le j\le d\)，记 \(\widehat H^{(j)}\) 为从 \(\widehat H\) 中删去第 \(j\) 列所得的 \(d\times d\) 子块，并定义最大子式综合征向量
\[
\Delta_j:=(-1)^{d+j}\det\!\bigl(\widehat H^{(j)}\bigr),
\qquad
\Delta:=(\Delta_0,\dots,\Delta_d)^{\mathsf T}\in k^{d+1}.
\]
则成立：
\begin{enumerate}
  \item \(\Delta\neq 0\) 且 \(\ker(\widehat H)=k\cdot \Delta\)；
  \item \(\Delta_d=\det(H_0)\neq 0\)。令
  \[
  p_j:=\Delta_j/\Delta_d\quad(0\le j\le d-1),
  \qquad
  p_d:=1,
  \]
  则存在唯一的首一 \(d\) 阶常系数递推满足
  \[
  \sum_{j=0}^{d}p_j\,a_{n+j}=0\qquad(\forall\,n\ge 0).
  \]
  \item 令
  \[
  H_1:=\bigl(a_{i+j+1}\bigr)_{0\le i,j\le d-1}\in M_d(k),
  \qquad
  A:=H_0^{-1}H_1\in M_d(k),
  \]
  并令多项式
  \[
  P(\lambda):=\sum_{j=0}^{d}p_j\,\lambda^{j}\in k[\lambda].
  \]
  则 \(P(A)=0\)；因此 \(A\) 的最小湮灭多项式为 \(d\) 次并与递推阶一致。
\end{enumerate}
\end{theorem}
\begin{proof}
(1) 由于 \(H_0\) 可逆，\(\widehat H\) 的秩为 \(d\)，从而 \(\dim\ker(\widehat H)=1\)。
将 \(\widehat H\) 的列向量记为 \(v_0,\dots,v_d\in k^d\)。
由 Laplace 展开恒等式（即 \(\bigwedge^{d+1}k^d=0\) 的坐标表达）得
\[
\sum_{j=0}^{d}(-1)^j\,v_j\det(v_0,\dots,\widehat v_j,\dots,v_d)=0.
\]
按定义 \(\Delta_j=(-1)^{d}\cdot (-1)^j\det(v_0,\dots,\widehat v_j,\dots,v_d)\)，故 \(\widehat H\Delta=0\)。
由 \(\dim\ker(\widehat H)=1\) 得 \(\ker(\widehat H)=k\cdot\Delta\)。
又因 \(\widehat H\) 的秩为 \(d\)，至少存在一个 \(d\times d\) 最大子式非零，故 \(\Delta\neq 0\)。
\par
(2) 由构造 \(\widehat H^{(d)}=H_0\)，故 \(\Delta_d=\det(H_0)\neq 0\)。
由 \(\widehat H\Delta=0\) 的第 \(i\) 行（\(0\le i\le d-1\)）得到
\(
\sum_{j=0}^{d}\Delta_j a_{i+j}=0
\)，等价于 \(\sum_{j=0}^{d}p_j a_{i+j}=0\)。
将无限 Hankel 矩阵的第 \(j\) 列记为
\(
c_j:=(a_{i+j})_{i\ge 0}
\)。
因 \(H_0\) 可逆，\(\{c_0,\dots,c_{d-1}\}\) 线性无关，并且存在唯一 \(p_0,\dots,p_{d-1}\) 使
\[
c_d+\sum_{j=0}^{d-1}p_j\,c_j=0.
\]
对该列关系施加移位算子（去掉首项）并归纳得到
\(
c_{n+d}+\sum_{j=0}^{d-1}p_j c_{n+j}=0
\)
对一切 \(n\ge 0\) 成立。
取第 \(0\) 个坐标即得 \(\sum_{j=0}^{d}p_j a_{n+j}=0\)（\(\forall n\ge 0\)）。
若另有首一 \(d\) 阶递推系数 \(p'_j\) 满足同一恒等式，则得到两组不同列关系，使 \(\{c_0,\dots,c_{d-1}\}\) 线性相关，矛盾；故递推唯一。
\par
(3) 由 \(H_0A=H_1\)，对 \(0\le j\le d-2\) 有
\[
H_0(Ae_j)=H_1e_j=\text{\(\widehat H\) 的第 \(j+1\) 列}=H_0e_{j+1},
\]
从而 \(Ae_j=e_{j+1}\)。
另一方面，由 (2) 的列关系可得
\(
\widehat H^{(d)}p+ \text{第 \(d\) 列}=0
\)，即
\(
H_0(p_0,\dots,p_{d-1})^{\mathsf T}+(a_d,\dots,a_{2d-1})^{\mathsf T}=0
\)；
注意 \((a_d,\dots,a_{2d-1})^{\mathsf T}\) 恰为 \(H_1e_{d-1}\)，故
\[
H_0(p_0,\dots,p_{d-1})^{\mathsf T}+H_1e_{d-1}=0
\quad\Rightarrow\quad
Ae_{d-1}=-(p_0,\dots,p_{d-1})^{\mathsf T}.
\]
因此 \(A\) 为递推多项式 \(P(\lambda)=\lambda^d+p_{d-1}\lambda^{d-1}+\cdots+p_0\) 的伴随矩阵形式，直接计算得到 \(P(A)=0\)。
又由 (2) 递推阶为 \(d\)，可知 \(A\) 的最小湮灭多项式次数至少为 \(d\)；而 \(\deg P=d\) 且 \(P(A)=0\) 给出次数至多为 \(d\)，故次数恰为 \(d\)。证毕。
\end{proof}

\begin{corollary}[良素数下最小递推多项式的保持性]\label{cor:xi-hankel-minimal-recurrence-good-prime-preservation}
\leanverified{paper\_xi\_hankel\_minimal\_recurrence\_good\_prime\_preservation}
取 \(k=\QQ\) 并设 \((a_n)_{n\ge 0}\subset\ZZ\) 在 \(\QQ\) 上的 Hankel 秩为 \(d\)，且 \(H_0\in M_d(\ZZ)\) 在 \(\QQ\) 上可逆。
令 \(P(\lambda)=\sum_{j=0}^{d}p_j\lambda^j\in\QQ[\lambda]\) 为定理 \ref{thm:xi-hankel-maximal-minor-syndrome-normal-form-uniqueness} 给出的唯一首一递推多项式。
取素数 \(p\nmid \det(H_0)\)，并记 \(\bar a_n:=a_n\bmod p\in\FF_p\) 与 \(\bar P:=P\bmod p\in\FF_p[\lambda]\)。
则对一切 \(n\ge 0\) 有
\[
\sum_{j=0}^{d}\bar p_j\,\bar a_{n+j}=0,
\]
并且 \(\bar P\) 为序列 \((\bar a_n)\) 在 \(\FF_p\) 上的最小首一递推多项式（因而其阶仍为 \(d\)）。
\end{corollary}

\begin{proof}
由定理 \ref{thm:xi-hankel-maximal-minor-syndrome-normal-form-uniqueness}(2)，存在整数 \(\Delta_0,\dots,\Delta_d\in\ZZ\) 使得
\[
\sum_{j=0}^{d}\Delta_j a_{n+j}=0\qquad(\forall\,n\ge 0),
\qquad
\Delta_d=\det(H_0)\neq 0,
\]
且 \(p_j=\Delta_j/\Delta_d\)。
逐项模 \(p\) 约化得到
\[
\sum_{j=0}^{d}\overline{\Delta_j}\,\bar a_{n+j}=0\qquad(\forall\,n\ge 0).
\]
由于 \(p\nmid \Delta_d\)，在 \(\FF_p\) 中有 \(\overline{\Delta_d}\neq 0\)，可将上式除以 \(\overline{\Delta_d}\) 得到所示首一递推，且系数恰为 \(\bar p_j\)。
\par
另一方面，\(\bar H_0:=H_0\bmod p\) 可逆。
令 \(\bar c_j:=(\bar a_{i+j})_{i\ge 0}\) 为 \(\FF_p\) 上无限 Hankel 矩阵的第 \(j\) 列。
若存在阶 \(r<d\) 的首一递推消去 \((\bar a_n)\)，则存在 \(u_0,\dots,u_{r-1}\in\FF_p\) 使得
\[
\bar c_r+\sum_{j=0}^{r-1}u_j\,\bar c_j=0,
\]
从而 \(\bar c_0,\dots,\bar c_{d-1}\) 线性相关，进而 \(\bar H_0\) 奇异，矛盾。
故最小阶必为 \(d\)，从而 \(\bar P\) 为最小首一递推多项式。证毕。
\end{proof}

\begin{corollary}[极小递推的 Hankel 主块可逆性与反演系数公式]\label{cor:xi-hankel-minimal-recurrence-inversion-formula}
\leanverified{paper\_xi\_hankel\_minimal\_recurrence\_inversion\_formula}
设 \(k\) 为域，\((a_m)_{m\ge 0}\subset k\) 满足某个\emph{极小阶}常系数线性递推：存在 \(d\ge 1\) 与系数 \(c_0,\dots,c_{d-1}\in k\) 使
\[
a_{m+d}=c_{d-1}a_{m+d-1}+\cdots+c_0a_m\qquad(m\ge 0),
\]
并且 \(d\) 在所有满足该性质的阶数中最小。
定义 Hankel 主块与右端向量
\[
H_d:=\bigl(a_{i+j}\bigr)_{0\le i,j\le d-1}\in M_d(k),
\qquad
b:=(a_d,a_{d+1},\dots,a_{2d-1})^{\mathsf T}\in k^d.
\]
则 \(H_d\) 必可逆，且系数向量 \(c:=(c_0,\dots,c_{d-1})^{\mathsf T}\) 由有限样本唯一确定并满足
\[
\boxed{\;c=H_d^{-1}b.\;}
\]
\end{corollary}
\begin{proof}
对 \(m=0,1,\dots,d-1\) 写出递推式得到线性系统
\(
H_d\,c=b
\)。
若 \(H_d\) 奇异，则存在 \(0\neq \alpha\in k^d\) 使 \(H_d\alpha=0\)，即
\(
\sum_{j=0}^{d-1}\alpha_j a_{i+j}=0
\)
对 \(0\le i\le d-1\) 成立。
由递推闭包可推出该关系对一切 \(i\ge 0\) 成立，从而给出阶 \(\le d-1\) 的非平凡递推，与极小性矛盾。
故 \(H_d\) 可逆，且 \(c\) 唯一并由 \(c=H_d^{-1}b\) 给出。证毕。
\end{proof}

\begin{theorem}[有限窗 Hankel 零空间的截断理想结构]\label{thm:xi-hankel-nullspace-truncated-ideal-structure}
\leanverified{paper\_xi\_hankel\_nullspace\_truncated\_ideal\_structure}
设 \(k\) 为域，\((a_n)_{n\ge 0}\subset k\) 的 Hankel 秩为 \(d<\infty\)。
假设存在某个偏移 \(\ell\ge 0\) 使得
\[
H_0^{(\ell)}:=\bigl(a_{\ell+i+j}\bigr)_{0\le i,j\le d-1}\in M_d(k)
\]
可逆。
令
\[
P(\lambda)=\lambda^{d}+c_{d-1}\lambda^{d-1}+\cdots+c_0\in k[\lambda]
\]
为 \(a\) 的首一极小递推多项式（等价地，对一切 \(n\ge 0\) 有 \(\sum_{j=0}^{d}c_j a_{n+j}=0\)，其中 \(c_d=1\)，且 \(d\) 极小）。
对任意 \(N\ge d\)，令
\[
H^{(N)}_{\ell}:=\bigl(a_{\ell+i+j}\bigr)_{0\le i,j\le N-1}\in M_N(k)
\]
为长度 \(N\) 的 Hankel 主块。
则存在自然同构
\[
\ker H^{(N)}_{\ell}\ \cong\ (P)\big/(\lambda^{N})\ \subset\ k[\lambda]/(\lambda^{N}),
\]
其中右端按 \(k\)-向量空间结构理解，并且该同构可取为：
对任意 \(\deg Q\le N-d-1\)，将 \(Q(\lambda)\) 送到向量
\[
\mathrm{coef}\bigl(Q(\lambda)P(\lambda)\bmod \lambda^{N}\bigr)\in k^{N}.
\]
因此 \(\dim_k\ker H^{(N)}_{\ell}=N-d\)，并且每一个长度为 \(N\) 的 Hankel 零向量都唯一对应于某个截断倍多项式 \(Q(\lambda)P(\lambda)\bmod\lambda^{N}\)。
\end{theorem}
\begin{proof}
记线性泛函 \(L:k[\lambda]\to k\) 为 \(L(\lambda^{t})=a_{\ell+t}\)。
对 \(x=(x_0,\dots,x_{N-1})^{\mathsf T}\in k^N\) 定义多项式
\(
Q_x(\lambda):=\sum_{j=0}^{N-1}x_j\lambda^j
\)。
则
\[
(H^{(N)}_{\ell}x)_i=\sum_{j=0}^{N-1}a_{\ell+i+j}x_j=L(\lambda^{i}Q_x(\lambda))
\qquad(0\le i\le N-1),
\]
从而 \(x\in\ker H^{(N)}_{\ell}\) 等价于 \(L(\lambda^{i}Q_x)=0\) 对所有 \(0\le i\le N-1\) 成立。
\par
由极小递推 \(\sum_{j=0}^{d}c_j a_{n+j}=0\) 可知对一切 \(t\ge 0\) 有 \(L(\lambda^{t}P(\lambda))=0\)。
因此对任意 \(\deg Q\le N-d-1\) 与任意 \(0\le i\le N-1\)，有
\(
L(\lambda^{i}Q(\lambda)P(\lambda))=0
\)；
又因 \(L(\lambda^{n})\) 仅依赖于幂指数 \(n\)，故截断 \(Q(\lambda)P(\lambda)\bmod \lambda^{N}\) 的系数向量必落在 \(\ker H^{(N)}_{\ell}\)。
这给出线性映射
\[
k[\lambda]_{\le N-d-1}\longrightarrow \ker H^{(N)}_{\ell},
\qquad
Q\longmapsto \mathrm{coef}\bigl(QP\bmod \lambda^N\bigr).
\]
其像同构于 \((P)/(\lambda^N)\)，并且 \(\dim (P)/(\lambda^N)=N-d\)。
\par
另一方面，由 Hankel 秩为 \(d\) 知任意有限 Hankel 块秩不超过 \(d\)，而 \(H_0^{(\ell)}\) 可逆给出 \(\mathrm{rank}(H^{(N)}_{\ell})\ge d\)，故 \(\mathrm{rank}(H^{(N)}_{\ell})=d\) 并得到
\(
\dim\ker H^{(N)}_{\ell}=N-d
\)。
维数比较迫使上述映射为同构，因而得到所述唯一性。证毕。
\end{proof}

\leanverified{paper\_xi\_hankel\_maximal\_minor\_syndrome\_offset\_projective\_invariance}
\begin{theorem}[最大子式综合征的偏移射影不变性与规范化唯一性]\label{thm:xi-hankel-maximal-minor-syndrome-offset-projective-invariance}
设 \(k\) 为域，\((a_n)_{n\ge 0}\subset k\) 的 Hankel 秩为 \(d<\infty\)。
对任意偏移 \(\ell\ge 0\)，定义移位 Hankel 主块与扩展块
\[
H_0^{(\ell)}:=\bigl(a_{\ell+i+j}\bigr)_{0\le i,j\le d-1}\in M_d(k),
\qquad
\widehat H^{(\ell)}:=\bigl(a_{\ell+i+j}\bigr)_{\substack{0\le i\le d-1\\ 0\le j\le d}}\in M_{d,d+1}(k).
\]
当 \(\det(H_0^{(\ell)})\neq 0\) 时，对 \(0\le j\le d\) 记 \(\widehat H^{(\ell,j)}\) 为从 \(\widehat H^{(\ell)}\) 中删去第 \(j\) 列所得的 \(d\times d\) 子块，并定义综合征向量
\[
\Delta_j^{(\ell)}:=(-1)^{d+j}\det\!\bigl(\widehat H^{(\ell,j)}\bigr),
\qquad
\Delta^{(\ell)}:=(\Delta_0^{(\ell)},\dots,\Delta_d^{(\ell)})^{\mathsf T}\in k^{d+1}.
\]
则存在与 \(\ell\) 无关的唯一向量 \(p=(p_0,\dots,p_{d-1},1)\in k^{d+1}\)，使得对一切满足 \(\det(H_0^{(\ell)})\neq 0\) 的 \(\ell\) 都有
\[
\boxed{\;\Delta^{(\ell)}=\det\!\bigl(H_0^{(\ell)}\bigr)\cdot p.\;}
\]
因此射影类 \([\Delta^{(\ell)}]\in\mathbf P^d(k)\) 与偏移 \(\ell\) 无关；等价地，由最大子式综合征抽取的首一 \(d\) 阶递推多项式在所有可逆偏移处完全一致。
\end{theorem}
\begin{proof}
固定满足 \(\det(H_0^{(\ell)})\neq 0\) 的 \(\ell\)，并考虑移位序列 \(a^{(\ell)}_n:=a_{\ell+n}\)。
此时 \(a^{(\ell)}\) 的 Hankel 秩仍为 \(d\)，且其对应主块正是 \(H_0^{(\ell)}\)，故可对 \(a^{(\ell)}\) 应用定理 \ref{thm:xi-hankel-maximal-minor-syndrome-normal-form-uniqueness}：
向量 \(\Delta^{(\ell)}\) 生成 \(\ker(\widehat H^{(\ell)})\)，并给出唯一首一递推系数
\[
p_j^{(\ell)}:=\frac{\Delta_j^{(\ell)}}{\Delta_d^{(\ell)}}\quad(0\le j\le d-1),
\qquad
p_d^{(\ell)}:=1,
\]
使得
\(
\sum_{j=0}^{d}p_j^{(\ell)}a_{\ell+n+j}=0
\)
对一切 \(n\ge 0\) 成立。
\par
现取任意两偏移 \(\ell_1,\ell_2\) 满足 \(\det(H_0^{(\ell_i)})\neq 0\)。
不妨设 \(\ell_2\ge \ell_1\)。
将 \(\ell_1\) 处递推在 \(n\mapsto n+(\ell_2-\ell_1)\) 后与 \(\ell_2\) 处递推相减，得到对一切 \(n\ge 0\) 成立的关系
\[
\sum_{j=0}^{d-1}\bigl(p_j^{(\ell_1)}-p_j^{(\ell_2)}\bigr)\,a_{\ell_2+n+j}=0.
\]
若 \(p^{(\ell_1)}\neq p^{(\ell_2)}\)，则右端给出一个阶 \(\le d-1\) 的非平凡常系数递推，对移位序列 \((a_{\ell_2+n})_{n\ge 0}\) 处处成立；
这与 \(\det(H_0^{(\ell_2)})\neq 0\) 所强制的 Hankel 秩为 \(d\)（从而最小递推阶为 \(d\)）矛盾。
故 \(p^{(\ell)}\) 必与 \(\ell\) 无关，记其公共值为 \(p\)。
\par
最后注意 \(\widehat H^{(\ell,d)}=H_0^{(\ell)}\)，故 \(\Delta_d^{(\ell)}=\det(H_0^{(\ell)})\neq 0\)，并且按定义
\(
\Delta^{(\ell)}=\Delta_d^{(\ell)}\cdot(p_0,\dots,p_{d-1},1)^{\mathsf T}
\)。
代入 \(\Delta_d^{(\ell)}=\det(H_0^{(\ell)})\) 即得盒中恒等式。证毕。
\end{proof}

\begin{corollary}[最小递推系数的共同分母整除所有满秩 Hankel 行列式]\label{cor:xi-hankel-recurrence-common-denominator-divides-determinants}
\leanverified{paper\_xi\_hankel\_recurrence\_common\_denominator\_divides\_determinants}
取 \(k=\QQ\) 并设 \((a_n)_{n\ge 0}\subset\ZZ\) 在 \(\QQ\) 上 Hankel 秩为 \(d\)。
令
\[
P(\lambda)=\lambda^d+c_{d-1}\lambda^{d-1}+\cdots+c_0\in\QQ[\lambda]
\]
为定理 \ref{thm:xi-hankel-maximal-minor-syndrome-normal-form-uniqueness} 抽取的唯一首一最小递推多项式。
令
\[
D:=\min\{N\in\NN:\ Nc_j\in\ZZ\ \forall\,0\le j\le d-1\}
\]
为系数共同分母。
则对任意满足 \(\det(H_0^{(\ell)})\neq 0\) 的偏移 \(\ell\ge 0\) 都有
\[
D\mid \det(H_0^{(\ell)}).
\]
因此 \(D\mid \gcd_{\ell:\det(H_0^{(\ell)})\neq 0}\det(H_0^{(\ell)})\)。
特别地，若存在 \(\ell\) 使 \(|\det(H_0^{(\ell)})|=1\)，则 \(P(\lambda)\in\ZZ[\lambda]\)。
\end{corollary}

\begin{proof}
由定理 \ref{thm:xi-hankel-maximal-minor-syndrome-offset-projective-invariance}，对任意满足 \(\det(H_0^{(\ell)})\neq 0\) 的 \(\ell\) 有
\[
c_j=p_j=\frac{\Delta_j^{(\ell)}}{\Delta_d^{(\ell)}}=\frac{\Delta_j^{(\ell)}}{\det(H_0^{(\ell)})}\qquad(0\le j\le d-1),
\]
其中 \(\Delta_j^{(\ell)}\in\ZZ\) 且 \(\Delta_d^{(\ell)}=\det(H_0^{(\ell)})\)。
因此每个 \(c_j\) 的分母都整除 \(\det(H_0^{(\ell)})\)，共同分母 \(D\) 亦整除 \(\det(H_0^{(\ell)})\)。
最后一条由 \(D\mid \pm 1\Rightarrow D=1\) 立得。证毕。
\end{proof}

\begin{theorem}[Hankel 奇异性的三步公开币完美零知识证明]\label{thm:xi-hankel-singularity-hvzk}
令 \(p\) 为素数，令 \(H\in \FF_p^{n\times n}\) 为 Hankel 矩阵，并记缺口
\[
\delta:=n-\mathrm{rank}(H)\ge 1.
\]
则存在一个三步公开币的 honest-verifier 完美零知识交互证明系统，用以证明语言谓词“\(\delta\ge 1\)”：
\begin{enumerate}
  \item \textbf{完备性：}接受概率至少为 \(1-p^{-\delta}\)。
  \item \textbf{健全性：}对任意不在语言内的输入（即 \(\mathrm{rank}(H)=n\)），单次交互的误受骗概率至多为 \(1/p\)。
  \item \textbf{零知识：}验证者视图可由不依赖任何核向量见证的模拟器精确生成（分布一致）。
\end{enumerate}
一种实现如下。验证者采样 \(u\sim \mathrm{Unif}(\FF_p^n)\) 并发送 \(u\)。
证明者在可行时取 \(s\in\ker(H)\) 使 \(u^{\mathsf T}s=1\)（若不存在则重启该轮），再采样 \(r\sim\mathrm{Unif}(\FF_p^n)\) 并发送承诺
\[
t:=(Hr,\ u^{\mathsf T}r)\in \FF_p^n\times \FF_p.
\]
验证者再采样挑战 \(e\sim\mathrm{Unif}(\FF_p)\)，证明者回应 \(z:=r+es\)。
验证者接受当且仅当
\[
Hz=Hr,\qquad u^{\mathsf T}z=u^{\mathsf T}r+e.
\]
\end{theorem}
\begin{proof}
设 \(\ker(H)\) 的维数为 \(\delta\)。
对随机 \(u\)，线性泛函 \(\ell_u:\ker(H)\to \FF_p\)、\(s\mapsto u^{\mathsf T}s\) 为零映射当且仅当 \(u\in(\ker(H))^\perp\)。
由于 \(\dim(\ker(H))^\perp=n-\delta\)，有
\(
\Pr[u\in(\ker(H))^\perp]=p^{-\delta}
\)。
因此除去概率 \(p^{-\delta}\) 的例外，存在 \(s\in\ker(H)\) 满足 \(u^{\mathsf T}s=1\)，此时检验恒等式
\[
Hz=H(r+es)=Hr,\qquad u^{\mathsf T}z=u^{\mathsf T}r+e
\]
成立，从而完备性为 \(1-p^{-\delta}\)。
\par
对健全性，设 \(\mathrm{rank}(H)=n\)，则 \(\ker(H)=0\)。
若对同一承诺 \(t\) 能对两条不同挑战 \(e\neq e'\) 给出通过检验的 \(z,z'\)，则由 \(Hz=Hr=Hz'\) 得 \(H(z-z')=0\)，并由第二条检验式得 \(u^{\mathsf T}(z-z')=e-e'\neq 0\)。
令 \(s:=(z-z')/(e-e')\)，则 \(Hs=0\) 且 \(u^{\mathsf T}s=1\)，矛盾。
因此对任意固定承诺 \(t\) 至多有一个挑战 \(e\) 可被接受，挑战均匀即单轮误受骗概率 \(\le 1/p\)。
\par
对完美零知识，构造模拟器：均匀采样 \((e,z)\in\FF_p\times\FF_p^n\)，并置
\[
t:=(Hz,\ u^{\mathsf T}z-e).
\]
则由构造恒有 \(Hz\) 与 \(u^{\mathsf T}z\) 通过检验。
在真实执行中，给定 \(e\) 时 \(z=r+es\) 随 \(r\) 均匀而在 \(\FF_p^n\) 上均匀；并且 \(t=(Hr,u^{\mathsf T}r)\) 满足 \(Hr=Hz\) 与 \(u^{\mathsf T}r=u^{\mathsf T}z-e\)。
因此真实分布与模拟分布完全一致，得到 HV 完美零知识。证毕。
\end{proof}

\begin{corollary}[常数轮与缺口线性轮数的不可兼容性]\label{cor:xi-hankel-hvzk-constant-round-vs-defect-linear-round}
在定理 \ref{thm:xi-hankel-singularity-hvzk} 的公开币 HVZK 模型下，“\(\mathrm{rank}(H)<n\)”语言存在常数轮交互证明，并且完备性失败概率随缺口 \(\delta\) 以 \(p^{-\delta}\) 衰减。
因此，任何在同一全信息 HVZK 口径下宣称最小交互轮数满足线性规律
\(
R_{\min}\asymp \delta
\)
的结论，都必须依赖额外的验证接口约束（例如：非交互口径、受限探测、流式一次读取或受限内存等），而不能由标准 HVZK 模型推出。
\end{corollary}

\begin{corollary}[整数 Hankel 缺口的圆维解释与最小补账轴数]\label{cor:xi-hankel-defect-cdim-minimal-ledger}
\leanverified{paper\_xi\_hankel\_defect\_cdim\_minimal\_ledger}
令 \(H\in\ZZ^{n\times n}\) 为整数 Hankel 矩阵，并视其为有限生成离散交换群同态
\[
f_H:\ZZ^n\to\ZZ^n,\qquad v\mapsto Hv.
\]
则其圆维缺陷满足
\[
\delta(f_H)=\cdim(\ker f_H)=n-\mathrm{rank}_{\QQ}(H).
\]
并且使 \((f_H,r):\ZZ^n\to\ZZ^n\oplus R\) 单射所需的最小补账轴数（即最小 \(\cdim(R)\)）恰等于该缺陷（定理 \ref{thm:xi-cdim-minimal-ledger-cost}）。
\end{corollary}
\begin{proof}
由 \(\cdim(A)=\dim_{\QQ}(A\otimes_{\ZZ}\QQ)\)（定义 \ref{def:circle-dimension} 的等价口径）与 \(\ker(f_H)\otimes\QQ=\ker(H:\QQ^n\to\QQ^n)\) 得
\[
\cdim(\ker f_H)=\dim_{\QQ}\ker(H:\QQ^n\to\QQ^n)=n-\mathrm{rank}_{\QQ}(H).
\]
最小补账轴数断言由定理 \ref{thm:xi-cdim-minimal-ledger-cost} 直接推出。证毕。
\end{proof}

\begin{theorem}[单素数提升的整数 Hankel 缺口 HVZK 证书与显式误接受上界]\label{thm:xi-hankel-singularity-one-prime-lift-hvzk}
\leanverified{paper\_xi\_hankel\_singularity\_one\_prime\_lift\_hvzk}
令 \(\widehat H\in\ZZ^{n\times n}\) 为整数 Hankel 矩阵，并定义语言
\[
\mathcal L_{\QQ}:=\bigl\{\widehat H:\ \mathrm{rank}_{\QQ}(\widehat H)<n\bigr\}.
\]
取一有限素数集合 \(\mathcal P\)，并令验证者公开抽样 \(p\sim\mathrm{Unif}(\mathcal P)\)。
令
\(
H_p:=\widehat H\bmod p\in\FF_p^{n\times n}
\)。
在给定 \(p\) 后，双方运行定理 \ref{thm:xi-hankel-singularity-hvzk} 的三步公开币 HVZK 协议以证明谓词 \(\mathrm{rank}(H_p)<n\)，并以其接受作为复合协议的接受判据。
则：
\begin{enumerate}
  \item 若 \(\widehat H\in\mathcal L_{\QQ}\)，则对任意 \(p\in\mathcal P\) 都有 \(\mathrm{rank}(H_p)<n\)，并且在条件于该 \(p\) 时复合协议接受概率至少为 \(1-p^{-\delta_p}\)，其中 \(\delta_p:=n-\mathrm{rank}(H_p)\ge 1\)。
  \item 若 \(\widehat H\notin\mathcal L_{\QQ}\)，令 \(\Delta:=\det(\widehat H)\in\ZZ\setminus\{0\}\)，则误接受概率满足
  \[
  \Pr[\text{误接受}]
  \ \le\
  \Pr_{p\sim\mathcal P}[\,p\mid \Delta\,]\ +\ \EE_{p\sim\mathcal P}\!\left[\frac{1}{p}\right].
  \]
  特别地，当 \(\mathcal P\) 为区间 \([B,2B]\) 内素数并作均匀抽样时，有
  \[
  \Pr[\text{误接受}]
  \ \le\
  \frac{\omega(\Delta)}{\#\mathcal P}\ +\ \frac{1}{B},
  \]
  其中 \(\omega(\Delta):=\#\{p:\ p\mid \Delta\}\) 为 \(\Delta\) 的不同素因子个数。
  \item 复合协议为 honest-verifier 完美零知识：存在模拟器先按同分布抽取 \(p\)，再调用定理 \ref{thm:xi-hankel-singularity-hvzk} 的模拟器生成交互记录；其分布与真实执行完全一致。
\end{enumerate}
\end{theorem}
\begin{proof}
若 \(\widehat H\in\mathcal L_{\QQ}\)，则 \(\det(\widehat H)=0\)，从而 \(\det(H_p)\equiv 0\pmod p\) 对任意 \(p\) 成立，故 \(\delta_p\ge 1\)。
条件于 \(p\) 时的接受率下界由定理 \ref{thm:xi-hankel-singularity-hvzk} 的完备性给出。
\par
若 \(\widehat H\notin\mathcal L_{\QQ}\)，则 \(\Delta\neq 0\)。
当 \(p\nmid \Delta\) 时，\(\det(H_p)\not\equiv 0\pmod p\)，从而 \(\mathrm{rank}(H_p)=n\)，内层命题为假；由定理 \ref{thm:xi-hankel-singularity-hvzk} 的健全性得条件接受概率 \(\le 1/p\)。
当 \(p\mid \Delta\) 时，内层命题为真，不能用健全性压制；对两类事件并合取上界即得
\[
\Pr[\text{误接受}]
\le \Pr_{p\sim\mathcal P}[\,p\mid \Delta\,]+\EE_{p\sim\mathcal P}[1/p].
\]
当 \(\mathcal P\) 为 \([B,2B]\) 内素数时，\(\EE[1/p]\le 1/B\)，且区间内满足 \(p\mid\Delta\) 的素数至多 \(\omega(\Delta)\) 个，故得到所示显式界。
\par
零知识部分由公开抽样 \(p\) 与内层 HVZK 的可模拟性直接组合：先抽 \(p\)，再模拟内层转录，分布保持一致。证毕。
\end{proof}

\begin{proposition}[\texorpdfstring{$p$}{p}-进赋值对模 \texorpdfstring{$p$}{p} 秩坍缩的数量控制]\label{prop:xi-det-vp-controls-rank-drop}
\leanverified{paper\_xi\_det\_vp\_controls\_rank\_drop}
设 \(A\in\ZZ^{d\times d}\) 满秩，并令 \(\Delta:=\det(A)\in\ZZ\setminus\{0\}\)。
对任意素数 \(p\) 记 \(A_p:=A\bmod p\in \FF_p^{d\times d}\)，并令
\[
\delta_p:=d-\mathrm{rank}_{\FF_p}(A_p).
\]
则有不等式
\[
\boxed{\;\delta_p\ \le\ v_p(\Delta)\;}
\]
并且当 \(p\nmid \Delta\) 时必有 \(\delta_p=0\)（即 \(A_p\) 在 \(\FF_p\) 上仍满秩）。
\end{proposition}
\begin{proof}
取 Smith 正规形：存在 \(U,V\in\mathrm{GL}_d(\ZZ)\) 使
\[
UAV=\mathrm{diag}(s_1,\dots,s_d),\qquad s_1\mid s_2\mid\cdots\mid s_d,
\]
且 \(\Delta=\prod_{i=1}^{d}s_i\)。
模 \(p\) 下，\(\mathrm{rank}_{\FF_p}(A_p)\) 等于对角元中不被 \(p\) 整除者的个数，因此秩亏损 \(\delta_p\) 等于满足 \(p\mid s_i\) 的指标个数。
每出现一次 \(p\mid s_i\) 至少贡献 \(v_p(s_i)\ge 1\)，从而
\[
\delta_p=\#\{i:\ p\mid s_i\}\le \sum_{i=1}^{d} v_p(s_i)=v_p(\Delta),
\]
得到所示不等式。
当 \(p\nmid\Delta\) 时 \(v_p(\Delta)=0\)，故 \(\delta_p=0\)。证毕。
\end{proof}

\begin{proposition}[\texorpdfstring{$p$}{p}-进求值的精度损失由 \texorpdfstring{$\det(H_0)$}{det(H0)} 精确控制]\label{prop:xi-hankel-padics-precision-loss-detH0}
\leanverified{paper\_xi\_hankel\_padics\_precision\_loss\_deth0}
令 \(p\) 为素数，取 \(d\ge 1\)。
设 \(H_0,H_1\in \ZZ_{p}^{d\times d}\) 且 \(\det(H_0)\neq 0\)，并令
\[
A:=H_0^{-1}H_1\in \QQ_{p}^{d\times d},
\qquad
s:=v_p(\det(H_0)).
\]
若仅给定 \(H_0,H_1\) 的模 \(p^N\) 截断（即在 \((\ZZ/p^N\ZZ)^{d\times d}\) 中的像），则在一般情形下最多只能确定
\[
A\pmod{p^{N-s}}.
\]
并且该精度损失 \(s\) 为尖锐：存在 \(d=1\) 的例子使得在给定同一模 \(p^N\) 数据条件下，\(A\) 的不确定性恰发生在 \(p^{N-s}\) 阶，无法提升到 \(p^{N-s+1}\)。
\end{proposition}
\begin{proof}
由伴随矩阵公式 \(H_0^{-1}=\mathrm{adj}(H_0)/\det(H_0)\) 得
\[
A=\frac{\mathrm{adj}(H_0)\,H_1}{\det(H_0)}.
\]
若 \(H_0,H_1\) 已知至模 \(p^N\)，则分子 \(\mathrm{adj}(H_0)H_1\) 亦至多确定到模 \(p^N\)。
除以 \(\det(H_0)=p^s\cdot u\)（\(u\in\ZZ_p^\times\)）使精度至多损失 \(s\) 位，从而 \(A\) 至多确定到模 \(p^{N-s}\)。
\par
尖锐性取 \(d=1\)：令 \(H_0=p^s\)，取 \(H_1=1\) 与 \(H_1'=1+p^N\)。
则 \(H_1\equiv H_1'\pmod{p^N}\)，但
\[
A=\frac{1}{p^s},\qquad
A'=\frac{1+p^N}{p^s}=A+p^{N-s},
\]
故不确定性不可消去于 \(p^{N-s}\) 阶。证毕。
\end{proof}

\begin{theorem}[判别式控制的 \(p\)-进节点分辨率半损耗律]\label{thm:xi-padic-root-resolution-discriminant-half-loss}
沿用命题 \ref{prop:xi-hankel-padics-precision-loss-detH0} 的设定，并令 \(M:=N-s\)。
设 \(K/\QQ_p\) 为某个分裂扩张，\(\mathcal O_K\) 为其赋值环，\(\pi\) 为均值元，赋值记为 \(v=v_\pi\)。
令 \(\chi(T)\in\mathcal O_K[T]\) 为某个首一 \(d\) 次多项式，并假设其在 \(K\) 中分裂为互异根 \(\{z_j\}_{j=1}^{d}\subset\mathcal O_K\)。
记判别式赋值
\[
D:=v\!\bigl(\operatorname{Disc}(\chi)\bigr).
\]
若仅给定 \(\chi\) 的系数向量在模 \(\pi^{M}\) 意义下的同余类，则根多重集 \(\{z_j\}\) 在最坏情形下至多可唯一确定到模 \(\pi^{M-\lceil D/2\rceil}\)，并且一般不能提升到模 \(\pi^{M-\lceil D/2\rceil+1}\)。
在 Gram--shift 实现口径下，将 \(\chi\) 取为 \(A\) 的特征多项式即可得到：在模 \(p^{N}\) 数据之下，节点谱的额外损耗由判别式赋值的半量纲 \(\lceil D/2\rceil\) 强制给出。
\end{theorem}

\begin{proof}
把根到系数的 Vieta 映射写为
\[
\Psi:\ \mathcal O_K^{d}\longrightarrow \mathcal O_K^{d},\qquad
(z_1,\dots,z_d)\longmapsto (c_1,\dots,c_d),
\]
其中 \(\chi(T)=T^{d}+c_1T^{d-1}+\cdots+c_d\)。
在互异根处，\(\det(D\Psi)\) 等于 Vandermonde 行列式 \(\prod_{i<j}(z_j-z_i)\) 的单位倍，因此
\[
v\!\bigl(\det(D\Psi)\bigr)=\sum_{1\le i<j\le d}v(z_j-z_i)=\frac12\,v\!\Bigl(\prod_{i<j}(z_j-z_i)^2\Bigr)=\frac{D}{2}.
\]
因此当系数仅给到模 \(\pi^{M}\) 时，任何试图反演 \(\Psi\) 的过程在最坏情形下都无法在 \(\pi\)-进意义下获得超过 \(M-\lceil D/2\rceil\) 的根分辨率；这是非阿基米德逆函数定理/多变量 Hensel 线性化中由雅可比赋值强制出现的精度损耗量纲。
\par
尖锐性可在二次情形显式给出。
取任意 \(e\ge 1\) 并令
\(
\chi(T)=(T-a)(T-a-\pi^{e})
=T^2-(2a+\pi^{e})T+a(a+\pi^{e})
\)
（\(a\in\mathcal O_K\) 任取），则 \(\operatorname{Disc}(\chi)=\pi^{2e}\) 从而 \(D=2e\)。
对任意满足 \(M\ge 2e+1\) 的整数 \(M\)，令 \(\delta:=\pi^{M-e}\) 并定义另一组互异根
\[
z_1':=a+\delta,\qquad z_2':=a+\pi^{e}-\delta.
\]
则 \(z_1'+z_2'=2a+\pi^e\) 与 \(z_1'z_2'=a(a+\pi^e)+\delta(\pi^e-\delta)\)。
由于 \(v(\delta\pi^e)=M\) 且 \(v(\delta^2)=2(M-e)\ge M+1\)，得
\(
z_1'z_2'\equiv a(a+\pi^e)\ (\mathrm{mod}\ \pi^{M})
\)，
从而以 \(\{z_1',z_2'\}\) 为根的首一二次多项式与 \(\chi\) 的系数向量在模 \(\pi^{M}\) 意义下完全一致。
另一方面，\(v(z_1'-a)=v(\delta)=M-e\)，故 \(z_1'\not\equiv a\ (\mathrm{mod}\ \pi^{M-e+1})\)（同理 \(z_2'\not\equiv a+\pi^e\ (\mathrm{mod}\ \pi^{M-e+1})\)），
从而根多重集不能由系数同余类唯一确定到模 \(\pi^{M-e+1}=\pi^{M-\lceil D/2\rceil+1}\)。
证毕。
\end{proof}

\begin{corollary}[多素数总预算下界由 \texorpdfstring{$\log|\det(H_0)|$}{log|det(H0)|} 完全刻画]\label{cor:xi-hankel-padics-global-budget-detH0}
设 \(H_0,H_1\in \ZZ^{d\times d}\) 且 \(\det(H_0)\neq 0\)，并令
\(
s_p:=v_p(\det(H_0))
\)。
对每个素数 \(p\) 设仅给定 \((H_0,H_1)\) 的模 \(p^{N_p}\) 截断数据。
若要求在所有坏素数通道上达到 \emph{非退化恢复}（即对一切 \(p\mid \det(H_0)\) 有 \(N_p>s_p\)，从而 \(A=H_0^{-1}H_1\) 在每个通道上至少可确定到模 \(p^{N_p-s_p}\) 的非平凡精度），则必要条件为
\[
\mathcal B\ >\ \mathcal S,
\qquad
\mathcal B:=\sum_{p} N_p\log p,
\qquad
\mathcal S:=\sum_{p} s_p\log p=\log|\det(H_0)|.
\]
\end{corollary}
\begin{proof}
对每个坏素数 \(p\mid \det(H_0)\)，命题 \ref{prop:xi-hankel-padics-precision-loss-detH0} 表明：若 \(N_p\le s_p\)，则通道精度在损失后不超过 \(p^{0}\) 级别，无法获得非平凡的 \(p\)-进确定性。
因此非退化恢复要求对所有 \(p\mid \det(H_0)\) 有 \(N_p>s_p\)。
两边乘以 \(\log p\) 并对所有素数求和，得到
\(
\sum_p N_p\log p>\sum_p s_p\log p
\)。
另一方面，由整数分解恒等式 \(|\det(H_0)|=\prod_p p^{v_p(\det(H_0))}=\prod_p p^{s_p}\)，取对数得 \(\sum_p s_p\log p=\log|\det(H_0)|\)。
合并即得所示必要条件。证毕。
\end{proof}

\begin{theorem}[低于 \texorpdfstring{$2d$}{2d} 窗口的完美零知识不可判别性]\label{thm:xi-hankel-window-2dminus1-perfect-zk-indistinguishability}
\leanverified{paper\_xi\_hankel\_window\_2dminus1\_perfect\_zk\_indistinguishability}
设 \(k\) 为域，取 \(b\in k^{2d-1}\) 满足 \(\det H_0(b)\neq 0\)，并令
\[
\mathcal F_b:=\bigl\{a\in k^{\NN}:\ \mathrm{rankH}(a)=d,\ \tau_{2d-1}(a)=b\bigr\}.
\]
设一族转录分布 \(\{P_a\}_{a\in\mathcal F_b}\) 给出某一审计协议的输出（验证器为对转录的可测判定函数）。
若该协议相对于 \(\pi:=\tau_{2d-1}\) 为完美零知识，即存在某个仅依赖 \(b\) 的分布 \(S_b\) 使得
\[
P_a=S_b\qquad(\forall\,a\in\mathcal F_b),
\]
则对任意非平凡子集 \(\varnothing\subsetneq \mathcal Q\subsetneq \mathcal F_b\)，不存在验证器 \(V\) 使得
\[
\Pr[V\ \mathrm{accept}\mid a]=1\ \ (\forall\,a\in\mathcal Q),
\qquad
\Pr[V\ \mathrm{accept}\mid a]=0\ \ (\forall\,a\in\mathcal F_b\setminus\mathcal Q).
\]
\end{theorem}
\begin{proof}
由命题 \ref{prop:xi-hankel-window-fiber-1d}，\(\mathcal F_b\) 至少含两条不同序列。
令 \(V\) 为任意验证器（可测函数），其接受概率为
\(
\Pr[V\ \mathrm{accept}\mid a]=\int V(t)\,\dd P_a(t)
\)。
完美零知识假设给出 \(P_a=S_b\) 对所有 \(a\in\mathcal F_b\) 成立，故上式在 \(\mathcal F_b\) 上恒等于 \(\int V\,\dd S_b\)，与 \(a\) 无关。
因此 \(V\) 不可能在同一纤维 \(\mathcal F_b\) 上对非平凡 \(\mathcal Q\) 给出全对全错判别。证毕。
\end{proof}

\begin{theorem}[纤维内误差恢复所需的信息泄漏下界]\label{thm:xi-hankel-fiber-fano-information-leakage}
\leanverified{paper\_xi\_hankel\_fiber\_fano\_information\_leakage}
令 \(k=\FF_Q\) 为有限域，取 \(b\in k^{2d-1}\) 满足 \(\det H_0(b)\neq 0\)。
由命题 \ref{prop:xi-hankel-window-fiber-1d}，存在自然双射 \(x\mapsto a^{(x)}\) 使得
\(\mathcal F_b=\{a^{(x)}:\ x\in\FF_Q\}\)。
令 \(X\sim\mathrm{Unif}(\FF_Q)\)，并令随机转录 \(T\) 由协议在输入 \(a^{(X)}\) 时生成。
若存在估计器 \(\widehat X=\widehat X(T)\) 满足 \(\Pr[\widehat X\neq X]\le \varepsilon\)（\(0<\varepsilon<1\)），则互信息满足
\[
\mathbf I(X;T)\ \ge\ \log Q\ -\ h(\varepsilon)\ -\ \varepsilon\log(Q-1),
\]
其中 \(h(\varepsilon):=-\varepsilon\log\varepsilon-(1-\varepsilon)\log(1-\varepsilon)\) 为二元熵函数（与 \(\log\) 取同一底）。
\end{theorem}
\begin{proof}
Fano 不等式给出
\[
H(X\mid T)\le h(\varepsilon)+\varepsilon\log(Q-1),
\]
且 \(H(X)=\log Q\)。
故
\[
\mathbf I(X;T)=H(X)-H(X\mid T)\ge \log Q-h(\varepsilon)-\varepsilon\log(Q-1).
\]
证毕。
\end{proof}
